\cotefile{29}

\scanpage{13}
1. Données de ramification sur un schéma.

Soit $S$ un schéma (préschéma dans l'ancienne terminologie), $\underline{\mathrm{Et}}_S$ le site étale de $S$. Nous appel\unsure{\deleted{r}e|\deleted{e}r}rons \uline{donnée de ramification sur} $S$ \deleted{la donnée} un couple $\inserted{\underline{R}=}(\underline{C},\underline{s})$, où $\underline{C}$ est une catégorie fibrée sur $\underline{\mathrm{Et}}_S$ (\deleted{\unsure{en particulier}}\inserted{i.e.} c'est une catégorie $\underline{C}$ munie d'un foncteur fibrant $p:\underline{C}\quad\underline{\mathrm{Et}}_S$) et où $\underline{s}:\underline{C}\quad\underline{\mathrm{Fl}}_S$ est un foncteur cartésien, où $\underline{\mathrm{Fl}}_S$ désigne la catégorie fibrée des \deleted{\unsure{flêches|flèches}} schémas relatifs sur les objets $S'$ de $\underline{\mathrm{Et}}_S$. On fera attention qu'on ne suppose pas que la catégorie fibrée $\underline{C}$ soit un champ (i.e. qu'on puisse recoller morphismes et objects de $\underline{C}$). Un objet $\inserted{M}$ de $\deleted{\underline{C}(S')}$ la catégorie fibre $\underline{C}_{S'}$ ($S'\quad\mathrm{Ob}\,\underline{\mathrm{Et}}_S$) sera appelé un modèle de ramification pour $\underline{R}$, et le schéma $\underline{s}(M)$ sur $S'$ sera appelé la réalisation géométrique de ce modèle. On fera attention qu'on ne suppose pas le foncteur $\underline{s}$ fidèle ni pleinement fidèle, donc les morphismes entre modèles de ramification ne peuvent être identifiés aux morphismes entre leurs réalisations géométriques. Dansles exemples envisagés par la suite, les $\underline{s}(M)$ sont finis sur $S'$.

Considérons le foncteur composé $\underline{C}\overset{\underline{s}}{\qquad}\underline{\mathrm{Fl}}_S\overset{\mathrm{source}}{\qquad}(\mathrm{Sch})$, d'où un foncteur (le foncteur ``réalisation géométrique'')
\[
\underline{C}\qquad(\mathrm{Sch})\,.
\]
Considérons d'autre part la catégorie fibrée sur$(\mathrm{Sch})$ des schémas relatifs étales et finis (dont la catégorie fibre en $T$ est donc la catégorie des revêtements étales finis de $T$), et soit $\underline{D}$ la catégorie fibrée sur $\underline{C}$ image inverse de la catégorie fibrée précédente par le foncteur réal'sation géométrique. La fibre de cette catégorie fibrée en l'objet $M$ de $\underline{C}$ est donc la catégorie des revêtements étales finis de la réalisation géométrique $\underline{s}(M)$ de $M$. O\textsuperscript{\unsure{N|n}} appelera \uline{revêtement de} \deleted{ty\unsure{p}}
% archive sheet number: 12

\scanpage{14}
Appelons \uline{crible couvrant} de $\underline{C}$ toute sous-catégorie \inserted{\unsure{\deleted{n}p|\deleted{p}n}\unsure{\deleted{e}l|\deleted{l}e|\deleted{d}l|\deleted{l}d}eine} $\underline{C}_o$ de $\underline{C}$ qui est un crible (i.e. pour toute flèche, l'origine appartient à $\underline{C}_o$ s'il en est de même de l'extrémité) et tel que les $\underline{p}(M)$ pour $M\quad\mathrm{Ob}\,\underline{C}_o$ recouvrent l'objet final $S$ de $\underline{\mathrm{Et}}_S$. Pour tout $\underline{C}_o$, soit $\mathrm{Rev}^{\underline{R}}_{\underline{C}_o}(S)$ la catégorie des sections cartésiennes de $\underline{D}$ au d essus de $\underline{C}_o$. Pour $\underline{C}_o$ var\unsure{\deleted{s}i|\deleted{i}s}able\deleted{\unsure{s}}, ces catégories forment \deleted{\unsure{un|on|ou} les fibres d'une}\inserted{un systême inductif de} catégories\deleted{fibrée\unsure{s} sur la catégorie ordonnée de ces $\underline{C}_o$} (NB si $\underline{C}'_o\quad\underline{C}_o$, on obtient un foncteur restriction évident $\mathrm{Rev}^{\underline{R}}_{\underline{C}_o}(S)\quad\mathrm{Rev}^{\underline{R}}_{\underline{C}'_o}(S)$), \deleted{et} dont la limite inductive est notée $\mathrm{Rev}^{\underline{R}}(S)$ et appelée \uline{catégorie des revêtements de} $S$ \uline{de type} $\underline{R}$.

Signaons que dans tous les cas que nous aurons à utiliser, les foncteurs de transition $\mathrm{Rev}^{\underline{R}}_{\underline{C}_o}(S)\quad\mathrm{Rev}^{\underline{R}}_{\underline{C}'_o}(S)$ sont pleinement fidèles, donc aussi les foncteurs $\mathrm{Rev}^{\underline{R}}_{\underline{C}_o}(S)\quad\mathrm{Rev}^{\underline{R}}(S)$, qui permettent donc d'identifier (à équivalence près) les $\mathrm{Rev}^{\underline{R}}_{\underline{C}_o}(S)$ à des sous-catégories pleines de la catégorie $\mathrm{Rev}^{\underline{R}}(S)$ des revêtements detype $S$. Nous supposerons \deleted{de plus} que pour tout $\underline{C}_o$, \deleted{le foncteur composé} et tout $X\quad\mathrm{Ob}\,\mathrm{Rev}^{\underline{R}}_{\underline{C}_o}(S)$, le foncteur composé $\underline{C}_o\overset{X}{\qquad}\underline{D}\overset{\mathrm{source}}{\qquad}(\mathrm{Sch})_{/S}$, (où le deuxième foncteur est le foncteur déduit du foncteur source sur la catégorie des s chémas finis étales relatifs,) admet une limite inductive dans $(\mathrm{Sch})_{/S}$, que nous noterons $\unsure{\deleted{r}\underline{r}|\underline{r}|\deleted{f}r|\deleted{r}f}(X)$, et que pour $C'_o\quad\underline{C}_o$, désignant par $X'$ la restriction de $X$ à $\underline{C}'_o$, l'homomorphisme canonique $r(X'_o)\quad r(X)$ est un isomorphisme. (Ces conditions seront satisfaites dans les exemples \inserted{(de généralité croissante)} 1 à 3 ci-dessous). Les différents foncteurs $r$ se r accorde\unsure{nt|\deleted{m}nt|\deleted{nt}m} alors et
% archive sheet number: 13

\scanpage{15}
\deleted{\uline{type} $\underline{R}$ \uline{de} $S$ une section cartésienne de $\underline{D}$ sur $\underline{C}$. Ces revêtements de type $\underline{R}$ forment une catégorie \inserted{que nous noterons $\mathrm{Rev}_{\underline{R}}(S)$.} \handwritten{(} notée suivant les notations standard $\underset{\handwritten{\longleftarrow}}{\mathrm{Lim}}\,\underline{D}/\underline{C}$\handwritten{)}\handwritten{\unsure{,|)}}. On définira la \uline{réalisation géométrique} d'un revêtement $\handwritten{\unsure{\underline{\varnothing}|\underline{\phi}|\underline{\Phi}}}$ de type $\underline{R}$ comme étant la limite inductive (si elle existe) dans $(\mathrm{Sch}$ du foncteur composé $\underline{C}\overset{\handwritten{\underline{\varnothing}}}{\qquad}\underline{D}\qquad(\mathrm{Sch})_{/S}$, où le deuxième de ces foncteurs est le foncteur source \deleted{\handwritten{(}associant à tout \unsure{revêtement}}(dé\inserted{d}uit du foncteur source sur la catégorie des schémas finis étales relatifs). Si pour tout $\handwritten{\underline{\varnothing}}$, cette limite inductive existe,} on obtient un foncteur réalisation géométrique
\[
\deleted{\handwritten{\unsure{\underline{s}|\underline{\xi}|\phi|r}}}\quad r:\mathrm{Rev}_{\underline{R}}(S)\qquad(\mathrm{Sch})_{/S}\,.
\]
Dans les cas qui nous intéresseront par la suite, ce foncteur $r$ sera fidèle, mais pas nécessairement pleinement fidèle; il importera donc de bien distinguer entre une revêtement de type $\underline{R}$ de $S$, et sa réalisation géométrique, le premier n'étant pas déterminée par la seule connaissance du second.

Avant d'aller plus loin, donnons quelques exemples.

\uline{Exemple} 1. Donnons-nous un schéma relatif $Z$ sur $S$, muni d'un groupe fini $G$ opérant \inserted{(à droite disons)} par $S$-automorphismes. Soit $C_o$ la catégorie réduite à un objet $\inserted{e}$ et dont le monoïde des endomorphismes est $G$, et prenons pour $\underline{C}$ la catégorie produit $\underline{\mathrm{Et}}_S\mathrm{x}C_o$, \deleted{pour foncteur $\underline{s}$ le foncteur $(S',e)$ qui \unsure{sur|aux} les objets} La donnée d'un foncteur cartésien de \deleted{Et} $\underline{C}$ dans une catégorie fibrée $\underline{C}'$ sur $\underline{\mathrm{Et}}_S$ équivaut manifestement à la donnée d'un objet de la catégorie fibre $\underline{C}'_S$ avec une opération de $G$ dessus, de sorte que $(Z,G)$ nous fo\unsure{\deleted{rm}ur|\deleted{ur}rm}nit bien un foncteur cartésien $\underline{s}$. On trouve aussitôt que la catégorie $\mathrm{Rev}_{\underline{R}}(S)$ est équivalente canoniquement à la \deleted{\unsure{mét|rét|cét}} catégorie des revêtements finis ét\unsure{\deleted{q}a|\deleted{a}q}les $Z'$ de $Z$, munis d'une opération (droite) de $G$ \deleted{comp\unsure{a}tible} rendant le morphisme structural $Z'\quad Z$ un $G$-morphisme.
% archive sheet number: 14

\scanpage{16}
La réalisation géométrique d'un tel $Z'$ existe si et seulement si le quotient $Z'/G$ existe, \inserted{et elle s'identifie alors à $X=Z'/G$.} Il en est ainsi si toute orbite de $G$ dans $Z$ est contenue dans un ouvert affine, ce qui implique en effet qu'il en est de même de toute orbite de $G$ dans $Z'$. \deleted{Intuitivement, \inserted{et en première approximation, les réalisations} la considération les revêtements de type $\underline{R}$ correspondent à des revêtements} Dans les cas que nous avons en vue, $Z$ est un revêtement de $S$, donc la condition sur les orbites sera automatiquement satisfaite (il suffit que $Z$ soit affine sur $S$)\deleted{plus généralement que \unsure{d|s}es fibre\unsure{s}} de plus les réalisations géométriques $X=Z'/G$ seront automatiquement finies sur $S$, \deleted{donc}\inserted{i.e.} \unsure{d|\deleted{e}d|\deleted{d}e}es revêtements. En première approximation, ce seront les revêtements de $S$ qui (en \deleted{\unsure{a|u}jn} un sens convenable ...) deviennent étales au dessus du revêtement (ramifié en général) donné $Z$, ce qui donne un sens intuitif à la notion de ``donnée de ramification''. Soit maintenant $U$ un ouvert de $S$ tel que $Z\;U$ soit un revêtement principal de groupe $G$, et supposons que $Z\;U=V$ soit dense dans $Z$ (plus généralement, que toute partie ouverte et fermée de $Z$ contenant $Z\;U$ \inserted{et stable par $G$} soit égale à $Z$). Alors le foncteur $Z'\qquad Z'\;V$ de la catégorie des $G$-revêtements étales de $Z$ dans la catégorie des $G$-revétments étales de $V$ est fidèle, d'autre part le foncteur $Z'_V\qquad Z'_{\deleted{V}}/G$ de la catégorie des $G$-revêtement s étales de $V$ dans la catégorie des revêtements étales de $U$ est une équivalence de catégories, à fortiori fidèle, d'où on conclut que le foncteur \inserted{réalisation géométrique} $Z'\qquad Z'/G=X$ est\unsure{\deleted{f}|x} \uline{fidèle}. Suivant le yoga général\deleted{\unsure{e|m}}, \deleted{une donn\unsure{é|e}e de}\inserted{la catégorie des} revêtementsde type $\underline{R}$ \deleted{du} est alors équivalente à la catégorie des ``revêtements de $S$ munis d'une structure de type $\underline{R}$'', où par ``structure de type $\underline{R}$'' sur un revêtement $X$ de $S$ on entend la donnée d'une classe (mod isomorphisme) d'un $Z'\quad\mathrm{Rev}_{\underline{R}}(S)$ et d'un isomorphisme de $\underline{r}(Z')$ ($=Z'/G$) avec $X$. (Cette
% archive sheet number: 15

\scanpage{17}
dernière remarque aurait dû figurer avant l'exemple, elle s'applique chaque fois que le foncteur $\underline{r}$ est fidèle). \deleted{Nous pouvons également} Supposons maintenant $V$ dense dans $Z$ et de plus $Z$ \uline{normal}, alors on voit aussitôt que le foncteur $Z'\qquad Z'\;V$ est pleinement fidèle, d'où résult facilement qu'il en est de même du foncteur réalisation géométrique $Z'\qquad Z'/G$; donc dans ce cas, la catégorie $\mathrm{Rev}_{\underline{R}}(S)$ \deleted{\unsure{etd|età|est}} estéquivalente à une sous-catégorie pleine de la catégorie des revêtements \deleted{\unsure{de|du}} normaux de $S$,\inserted{étales au dessus de $U$, et} dont toute composante \deleted{domine} irréductible domine une composante irréductible de $S$ rencontrant $U$.

\uline{Exemple} 2 C'est une simple variante de l'exemple précédent: au lieu de se donner un $S$-schéma $Z$ à groupe fini d'opérateurs $G$, on se donne un systême projectif $(Z_i,G_i)_{i\;I}$ de tels objets ($I$ ensemble ordonné filtrant). \deleted{La catégorie Rev} Le systême projectif de\unsure{\deleted{s}|x|s|\deleted{s}x|\deleted{x}s} groupes $(G_i)$ définit une catégorie fibrée $C_o$ sur la catégorie associée à $I$ (les fibres de $C_o$ étant les catégories à un seul objet définies par les $G_i^o$), et on prendra $\underline{C}=\underline{\mathrm{Et}}_S\mathrm{x}C_o$, et on définit $\underline{s}$ de façon évidente. La catégorie $\mathrm{Rev}_{\underline{R}}(S)$ est maintenant équivalente à la catégorie $\mathrm{Lim}$ des catégories $\mathrm{Rev}_{\underline{R}_i}(S)$ définies séparément par les $(Z_i,G_i)$. Un cas intéressant est celui où les foncteurs transition (qui s'interprètent de façon évidente comme des foncteurs imagesinverses de revêtements étales à opérateurs) sont pleinement fidèles, de sorte que les $\mathrm{Rev}_{\underline{R}_i}(S)$ s'identifient (à équivalence près) à des sous-catégories pleines de $\mathrm{Rev}_{\underline{R}}(S)$ qui epuisent évidemment cette dernière. Il en sera ainsi si les morphismes de transition $(Z_i,G_i)\qquad(Z_j,G_j)$ sont tels que $G_i\quad G_j$ soit surjecti et $Z_i\quad Z_j$ induit un isomorphisme $Z_j\quad Z_i/N_{ij}\mathord{,}_{\unsure{R|\deleted{R}}}$ où $N_{ij}$ est le noyau de $G_i\quad G_j$ (en effet, il en résulte que pour tout revêtement \inserted{étale $Z'_j$}
% archive sheet number: 16

\scanpage{19}
de $Z_j$, on a $Z'_j\quad Z'_{i\deleted{j}}/N_{ij}$, où $Z'_i$ est l'image inverse de $Z'_{\unsure{\deleted{i}j|\deleted{j}i}}$ sur $Z_i$). Dans ce cas, le foncteur ``réalisation géométrique'' est bien défini, et induit surchaque $\mathrm{Rev}^{\underline{R}_i}(S)$ le foncteur réalisation géométrique déjà explicité dans l'exemple \unsure{\deleted{1}2|\deleted{2}1}. \deleted{Bien sûr, les conditions} Notons que les \unsure{cri\deleted{t}è\deleted{è}tres|cri\deleted{è}t\deleted{t}ères|cri\deleted{t}è\deleted{t}ères} signalés dans l'exemple 1 pour que le foncteur réalisation géométrique soit fidèle resp. pleinement fidèle, s'étendent de façon immédiate au cas envisagé ici.

\uline{Exemple} 3. Supposons que $\underline{R}=(\underline{C},\underline{s})$ satisfasse aux conditions suivantes. Le foncteur $\underline{s}$ se factorise en $\underline{C}\overset{\underline{s}'}{\qquad}\underline{\mathrm{F\unsure{\deleted{a}l|\deleted{l}a|\deleted{o}l|\deleted{l}o}op}}_S\qquad\underline{\mathrm{Fl}}_S$, où $\underline{\mathrm{Flop}}_S$ désigne la catégorie fibrée sur $\underline{\mathrm{Et}}_S$ des schémas relatifs à groupe fini (opérant à droite) d'opérateurs, et $\underline{\mathrm{Flop}}_S\quad\underline{\mathrm{Fl}}_S$ est le foncteur ``oubli\deleted{\unsure{r|e}} des opérations''. \deleted{Pour toute flèche $M'\quad M$ de $\underline{C}_{S'}$, la flèche \unsure{correspondante}}\inserted{Il faut préciser que (en vue de l'exemple kummérien plus} bas) nous prenons comme morphismes d'un schéma à opérateurs $(Z,G)$ sur $S$ dans un autre $(Z',G')$ les couples $(f,u)$, où $f:X\quad X'$ est un $S$-morphisme et $u:G_S\quad G'_S$ un morphisme(localement constant,\deleted{\unsure{(|f}}pas nécessairement constant !) de schémas en groupes sur $S$, de telle façon que $f$ soit compatible avec $u$ dans le s ens évident. On suppose que pour toute flèche $M'\quad M$ dans un $\underline{C}_S$, \deleted{déf} la flêche correspondante $\underline{s}(M')=(Z',G')\quad\underline{s}(M)=(Z,G)$ induit un isomorphisme $Z\quad Z'/H$, où $H$ est le noyau de $G'_S\quad G_S$ (qui est un sous-groupe localement constant de $G'_S$). On suppose \deleted{que tout} aussi que pour tout objet $M$ de $\mathrm{Ob}\,\underline{C}_{S'}$, posant $\underline{s}(M)=(Z,G)$, \deleted{et pour tout $g\quad G$, il existe un endomorphisme $\underline{g}$ de $M$, uniquement déterminé, tel que}\inserted{on se donne un \deleted{automorphisme} homomorphisme $G\quad\mathrm{Aut}(M)$, tel que pour tout $g\quad G$, si $\underline{g}$ est son image,} $\underline{s}'(\underline{g})$ soit le couple $(g_Z,\mathrm{int}(g))$, et \unsure{\inserted{tel} que tout|que tout \inserted{tel}} endomorphisme de $M$ soit localement de la forme $\underline{g}$; on \deleted{veut} suppose de plus que \deleted{la} l'homomorphisme $G\quad\mathrm{Aut}(M)$ soit compatib\unsure{\deleted{l}d|\deleted{d}l|\deleted{o}l|\deleted{l}o}e avec images inverses, et que pour toute flèche $M'\quad \unsure{M|\deleted{M}M|\deleted{X}M|\deleted{M}X}$ dans $\underline{C}_{S'}$, considérons $M'$ (resp. $M$) comme muni ud groupe $G'$ ($G$) d'opérateurs, $M'\quad M$ soit compatible avec $G'\quad G$.
% archive sheet number: 18

\scanpage{20}
Dans ce cas, le foncteur ``réalisation géométrique'' est bien défini et induit sur chaque $\mathrm{Rev}_{\underline{R}_i}(S)$ le foncteur réalisation géométrique déjà explicité dans \hspace{3em} l'exemple précédent.

de $Z_j$, on a $Z'_j\quad Z'_i/N_{ij}$, où $Z'_i$ est l'image inverse de $Z'_j$ sur $Z_i$). Notons que les critères signalés dans l'exemple 1 pour que le foncteur $\underline{r}:\mathrm{Rev}_{\underline{R}}(S)\qquad(\mathrm{Sch})_{/S}$ soit fidèle resp. pleinement fidèle, s'étendent de façon immédiate au cas envisagé ici.

\uline{Exemple} 3 Supposons que $\underline{R}=(\underline{C},\underline{s})$ satis\unsure{\deleted{g}f|\deleted{f}g}\unsure{\deleted{q}a|\deleted{a}q}sse aux conditions suivantes: pour tout $S'\quad\mathrm{Ob}\,\underline{\mathrm{Et}}_S$, et $M\quad\mathrm{Ob}\,\underline{C}_{S'}$, le groupe des automorphismes $G$ de $M$ est fini, et pour toute flèche $M'\quad M$ de $\underline{C}_{S'}$, il existe un unique morphisme de groupe $G'=\mathrm{Aut}(M')$ dans $G=\mathrm{Aut}(M)$ tel que $M'\quad M$ soit compatible avec les opérations de $G'$, de plus si $H$ est son noyau, le morphisme $\underline{s}(M')\quad\underline{\unsure{s|\deleted{e}s|\deleted{s}e}}(M)$ induit un isomorphisme $\underline{s}(M')/H\quad\underline{s}(M)$. \deleted{Enfin} On suppose \deleted{la catégorie $\underline{C}_{S'}$ filtrante (i.e. de pour deux objets $M$, $M'$ de ladite, il existe un objet $M''$ et des flèches $M\quad M''$, $M'\quad M''$),}\inserted{\deleted{que} de plus que pour deux morphismes de $M'$ dans $M^{\deleted{\prime}}$, \deleted{les morphismes correspondants de $G'$ dans $G$ sont \unsure{conjugués,|conjugués} déduit l'un de l'autre par composition avec un}} l'un est déduit de l'autre par composition avec un élément de $G$. Enfin, on suppose qu'étant donnés deux objets $M$; $M'$ de $\underline{C}_{S'}$, il existe une famille surjective de morphismes étales $S'_i\quad S'$ telle que pour tout $i$, $M_{S'}$ et $M'_{S_i}$ soient majorés par un objet $M''_i$ de $\underline{C}_{S'_i}$. Utilisant la théorie de la descente, on voit sous ces conditions que \unsure{\deleted{d}l|\deleted{l}d}\unsure{\deleted{e}a|\deleted{a}e} foncteur réalisation géométrique \deleted{de} $\mathrm{Rev}_{\underline{R}}(S)\qquad(\mathrm{Sch})_{/S}$ est défini, et à valeurs dans les \deleted{schémas fini\unsure{s} sur} revêtements de $S$ (en supposant que les $\underline{s}(M)$ sont des revêtements, bien sûr; pour pouvoir appliquer la descente, il suffirait même que les $\underline{s}(M)$ soient affines sur $S'$). Ici encore, on a des critères pour que ce foncteur soit fidèle resp. pleinement fidèle: il suffit que pour tout $S'\quad\mathrm{Ob}\,\underline{\mathrm{Et}}_S$ et tout $M\quad\mathrm{Ob}\,\underline{C}_{S'}$, la réalisation géométrique $\underline{s}(M)$, munie du groupe d'automorphismes $G=\mathrm{Aut}(M)$, satisfasse aux conditions \inserted{correspondantes} envisagées dans l'exemple 1 pour $(Z,G)$.

\uline{Cas particuliers du type kummérien}: a) Supposons données des sections
% archive sheet number: 19

\scanpage{21}
% archive sheet number: 20
\deleted{$(a_i)_{i\quad I}$ de $\underline{O}_S$, un} une famille $(D_i)_{i\quad I}$ de diviseurs (de Cartier) sur $S$. \unsure{Bour|Pour|\deleted{P}Bour|\deleted{B}Pour} tout $S'$ étale sur $S$, on désigne par $\underline{C}_{S'}$ la catégorie suivante: les objets sont les \deleted{couples triples} \inserted{quadruples $M=$} $(\underline{a},\underline{n},\inserted{G,}\underline{u})$, où $\underline{a}=(a_i)_{i\quad I}$ est une famille d'équations pour les $D_{iS'}$, $\unsure{\underline{n}=\deleted{-}(n_i)_{i\quad I}|\underline{n}=\deleted{(}-n_i)_{i\quad I}}$ une famille d'entiers $1$ premiers aux caractéristiques résiduelles de $S'$, \deleted{et $u:G\quad{}_{\underline{n}}$} $G$ un groupe fini, et $u:G_{S'}\quad{}_{\underline{n}S'}$ un isomorphisme de schémas en groupes sur $S'$. 2tant donné un deuxième tel quadruplet \inserted{$M'=$} $(\underline{a}',\underline{n}',G',u')$, \deleted{un morphisme du second dans le premier est $M'$ dans $M$} $\operatorname{Hom}_{S'}(M',M)$ est vide s\unsure{\deleted{i}a|\deleted{a}i}uf si $\underline{n}'\quad\underline{n}$, auquel cas c'est l'ensemble des $\underline{b}=(b_i)_{i\quad I}$ tels que $\underline{b}^{n}=\underline{a}'\underline{a}^{-1}$. La composition des morphis\unsure{m|n}es,et la définition des foncteurs changements de base pour une flêche $S''\quad S'$ dans $\underline{\mathrm{Et}}_S$, sont clairs. Le foncteur car\unsure{\deleted{é}t|\deleted{t}é}\unsure{\deleted{t}é|\deleted{é}t}sien $\underline{C}\quad(\mathrm{Sch})_{/S}$ est défini en associant à $M=(\underline{a},\underline{n},G,\underline{u})$ le revêtement Kummérien élémentaire $Z_{\underline{a}}^{\underline{n}}$, \deleted{et à une} \deleted{flê} considéré comme un revêtement à groupe d'opérateurs $G$, grâce à l'isomorphisme $G\quad{}_{\underline{n}}$. On définit de faç\deleted{\unsure{a|o|x}}on évidente la valeur de $\underline{r}$ sur les \unsure{\deleted{c}f|\deleted{f}c}lêches de $\underline{C}_{S'}$, et sa compatibilité avec les changements de base. On voit alors qu'on est sous les conditions de l'exemple 3. De plus, le foncteur $\underline{r}$ est fidèle, et il est pleinement fidèle lorsque $S$ est \deleted{noethérien aux points de} localement noethérien, et de plus \unsure{\deleted{j}n|\deleted{i}n|\deleted{n}j|\deleted{n}i}ormal en les points des $\operatorname{supp}D_i$, et si de plus les $D_i$ sont réduits, et toute \inserted{comp\unsure{o|c}sante irréductible d\deleted{e}'un $D_i$ est distincte de toute comp. d'un $D_j$ pour $j\ne i$}

b) Supposons donné un diviseur $D$ sur $S$, \deleted{une famille couvrante $\unsure{S\quad S|S'\quad S}$} et considérons le faisceau \unsure{\deleted{é}s|\deleted{s}é|ś}ur $\underline{\mathrm{Et}}_S$ associé au préfaisceau qui, à tout $S'$, associe l'ensemble des \deleted{diviseurs} ensembles \deleted{$\unsure{\underline{Q}|\underline{\Omega}|\underline{\delta}|\underline{D}|\underline{O}}$} \inserted{$(D_i)$} $\theta$ de diviseurs sur $S'$ ayant les propriétés suivantes: $\unsure{D_{S\overset{=}{'}}=\quad D_i|D_{S\deleted{=}'}=\quad D_i|D_{S\deleted{'}=} =\quad D_i|D_{S\bar{'}}=\quad D_i}$, et pour tout $s'\quad S'$, les restrictions \inserted{\deleted{non nulles}} à $\text{Sêc}(\underline{O}_{S',s'})$ des $D_i$ qui sont non nulles, sont deux à deux disjointes. Soit donnée une section du

\scanpage{22}
% archive sheet number: 21
faisceau associé au préfaisceau précédent (NB la condition canularesque mise sur les $(D_i)$ assure que l'on a bien un préfaisceau). Alors localement pour la topologie étale, on a donc une situation $(D_i)$ comme dans a). Cela suffit pour définir une donnée de ramification $\underline{R}$ sur $B$, généralisant d'ailleurs celle envisagée dans a) (du moins lorsque la famille $(D_i)_{i\quad I}$ de a) satisfait à la condition supplémentaire locale qu'on\handwritten{\unsure{g|?}} vient de formuler). Je vois deux cas particuliers ayant de l'intérêt;. \deleted{\unsure{S|R|fi|si}} et \deleted{1\textsuperscript{\uline{o}})} où $\quad$ peut être défini en termes d'un diviseur $D$ donné sur $S$:

1\textsuperscript{\uline{o}}) $S$ est localement noethérien, $D$ est à croisements normaux, on prend comme $D_i$ les composantes irréductibles locales régulières de $D$ (qui ne sont définies que localement pour la topologie étale).

2\textsuperscript{\uline{o}}) $S$ est au dessus d'un préschéma $T$, $D$ est à croisements normaux relativement à $T$, on prend comme $D_i$ les composantes locales lisses sur $\unsure{\deleted{T}S|\deleted{S}T}$ (qui ne sont définies que localement pour la topologie étale).

Dans le cas 1\textsuperscript{\uline{o}}), la catégorie $\operatorname{Rev}(_{\underline{R}}(S)$ est équivalente à une sous-catégorie pleine de la catégorie des revêtements de $S$, et sauf erreur (il faudrait le vérifier) il en est de même dans le cas 2\textsuperscript{\uline{o}}). Les revêtements en question s'appellent, dans le\deleted{ur} cas 1\textsuperscript{\uline{o}}), revêtements modérément ramifiés au sens large relativement au diviseur $D$ (NB on vérifie en effet que c'est une notion plus large que celle de ``revêtement modérément ramifié par rapport à $D$'', (signifiant : par rapport à la famille réduite à $D$). Le théorème d'Abhyankar permet de caractériser ces revêtements par la condition d'être étales sur \deleted{$S\quad\unsure{\underline{r}|\underline{a}|\underline{R}}$} $S-\operatorname{supp}D$, normaux en les points au dessus de $\operatorname{supp}D$, toute composante irréductible en d\unsure{\deleted{i}o|\deleted{o}i}mine une de $S$, enfin ils sont modérément ramifiés au sens classique en le\inserted{s} points maximaux de $D$. Dans le cas 2\textsuperscript{\uline{o}}, il y aurait lieu d'appeler les revêtements obtenus ``modérément ramifiés pour $D$, relativement à la

\scanpage{23}
% archive sheet number: 22
base $T$ ''. Lorsque $T$ est régulier, les notions 1\textsuperscript{\uline{o}} et 2\textsuperscript{\uline{o}} coïncident. Le théorème de spécialisation pour le groupe fondamental tame peut se démontrer, par la technique de SGA 1, quand on utilise la notion de revêtement modérément ramifié au sens large, relativement aux\unsure{diviseurs|diviseur} $D_{\bar{t}}$ \inserted{dans les $S_{\bar{t}}$} ($D$ étant à croisement normaux relativement à $T$, $S$ propre et lisse sur $T$).

\uline{Remarques} Evidemment, pour tout $S'$ objet de $\underline{\mathrm{Et}}_S$, on définit de façon évidente $\operatorname{Rev}_{\deleted{\underline{R}}}^{\inserted{\underline{R}}}(S')$, en considérant la restriction de $\underline{R}$ à $\underline{\mathrm{Et}}_{S'}\quad(\underline{\mathrm{Et}}_S)_{/S'}$. Pour $S'$ variable, on obtient ainsi une catégorie fibrée $\underline{\mathrm{Rev}}^{\underline{R}}$ sur $\underline{\mathrm{Et}}_S$, dont les fibres $\underline{\mathrm{Rev}}_{S'}^{\underline{R}}$ sont les $\operatorname{Rev}^{\underline{R}}(S')$, et un foncteur cartésien
\[
\underline{r}:\quad\underline{\mathrm{Rev}}^{\underline{R}}\qquad\underline{\mathrm{Rev}},
\tag*{\unsure{(*)|(\deleted{=}*)|(\deleted{*}=)}}
\]
où on désigne pour abréger par $\underline{\mathrm{Rev}}$ la catégorie \inserted{fibrée sur $\underline{\mathrm{Et}}_S$} des revêtements des différents objets $S'$ de $\underline{\mathrm{Et}}_S$. (NB on suppose que pour tout $S'$, le foncteur réalisation géométrique sur $\underline{C}_{S'}$ est à valeurs dans les revêtements de $S'$, et que le foncteur réalisation géométrique sur $\operatorname{Rev}^{\underline{R}}(S')$ est bien défini). Il faut insister sur le fait que pour nous, la ``donnée de ramification'' $\underline{R}$ est un objet de nature essentiellement auxiliaire, destiné à définir la catégorie fibrée $\underline{\mathrm{Rev}}^{\underline{R}}$ (qui est même un \uline{\unsure{c\deleted{h}a\deleted{a}hmp|c\deleted{a}h\deleted{h}amp}}) et le foncteur cartésien $\underline{r}$ de (*); c'est \unsure{(*)|(\deleted{=}*)|(\deleted{*}=)} qui est l'objet véritable de l'étude. On regardera donc deux données de ramification $\underline{R}$, $\underline{R}'$ comme essentiellement équivalentes, si elles définissent des foncteurs (*) équivalents (par une équivalence, bien sûr, induisant l'identité sur $\underline{\mathrm{Rev}}$). On peut se convaincre sur des exemples que $\underline{R}$ et $\underline{R}'$ peuvent avoir des allures très différentes. Ainsi, sous les conditions de l'exemple 3), si $\underline{R}'$ est défini par une sous-catégorie fibrée \deleted{de} $\underline{C}'$ de $\underline{C}$

\scanpage{24}
% archive sheet number: 23
telle que pour tout objet $M$ d'un $\underline{C}_{S'}$, il existe un recouvrement de $S'$ par des $S'_i$ tel que pour tout $i$, $M_{S'_i}$ soit \deleted{dominé} majo\unsure{r}é par un $M_i$ objet de $\underline{C}'_{S'_i}$, \deleted{- alors} et supposant que le foncteur réalisation géométrique pour $\underline{R}'$ est celui induit pour le foncteur réal. géom. pour $\underline{R}$, - alors on constate que $\underline{R}$ et $\underline{R}'$ sont ``équivalents''. Comme cas particulier, signalons l'exemple de la donnée kummérienne a), dans le cas où les $D_i$ admettent des équations globales par des sections $c_i$ de $\underline{O}_S$. Alors on peut considérer le systême projectif des $Z_{\underline{c}}^{\underline{n}}$, pour $\underline{n}=(n_i)_{i\quad I}$ variables, premiers \deleted{aux \unsure{caractéristiques|caractéritiques} résiduelles} à la caractéristique de $S$ (supposée exister). Considérons de plus le systême projectif des groupes $\underline{Z}_{\underline{n}}=\quad\underline{Z}/n_i\underline{Z}$, et supposons donné un isomorphisme entre celui-ci et le systême projectif des ${}_{\underline{n}}$ (il en existe lorsque $S$ est un schéma sur un corps algébriquement clos, par exemple). On peut alors considérer les $Z_{\underline{c}}^{\underline{n}}$ comme des schémas à groupes $\underline{Z}_{\underline{n}}$, formant un systême projectif, \deleted{d'où} donc on est sous les conditions de l'exemple 2, définissant une donnée de ramification $\underline{R}'$ sur $S$. On constat alors que cette donnée est équivalente à la donnée $\underline{R}$ décrite dans le cas particulier kummérien a). Comme autre exemple de données de ramification équivalentes\deleted{\unsure{:|;|.}}, signalons le suivant, où $\underline{R}$ est la donnée de ramification considérée dans l'exemple 1. On se donne un revêtement principal $P$ de $S$, de groupe $\unsure{H|\deleted{H}G|\deleted{G}H}'$ (opérant à droite), et on suppose que $G$ opère dessus librement à gauche, en commutant aux opérations de $G'$. Soit alors $Z'=Zx^{G}P$, considéré comme revêtement à groupe d'opérateurs $G'$, d'où une donnée de ramification $\underline{R}'$. Cette dernière est équivalente à $\underline{R}$, comme on voit en considérant l\deleted{es}'équivalences \deleted{quasi-inverses l'une} de \deleted{l'a\unsure{u}tre} \deleted{entre} $\underline{\mathrm{Rev}}^{\underline{R}}$ \deleted{et}\inserted{dans} ${\unsure{\underline{\mathrm{Rev}}|\mathrm{Rev}}}^{\unsure{R'|\underline{R}'}}$ opérateurs \illegible{3}

\scanpage{25}
% archive sheet number: 24
évidente lorsque $P=G'_S$ avec l'opération régulière droite de $G'$, et l'opération gauche de $G$ est donnée par un monomorphisme $G\quad G'$; on se ramène à ce cas particulier par localisation étale.

On peut se proposer de normaliser la notion de donnée de ramification, pour éliminer l'arbitraire dans le choix d'un $\underline{R}$. Mais alors \unsure{\deleted{o}il|\deleted{i}ol|\deleted{b}il|\deleted{d}il|\deleted{i}bl|\deleted{i}dl} semble qu'on soit infailliblement conduit à prendre comme catégorie $\underline{C}$ la catégorie $\underline{\mathrm{Rev}}^{\underline{R}}$ elle-même qu'on se propose \unsure{d\deleted{'}e|d\deleted{e}'} construire et d'étudier ! (Comparer avec la situation analogue pour les relations entre sites et topos, l'objet auxiliaire étant le site, l'objet intrinsèque de nature géométriquenétant le topos associé à ce site).

2. \uline{Propriétés de la catégorie} $\operatorname{Rev}^{\underline{R}}(S)$. \uline{Le groupe fondamental} $ {}^{\underline{R}}_1(S,\quad)$.

On constate immédiatement que la catégorie $\operatorname{Rev}^{\underline{R}}(S)$ hérite de toutes les propriétés d'ex\unsure{a|\deleted{e}a|\deleted{a}e}\unsure{s|c|\deleted{s}c|\deleted{c}s}titude \deleted{déjà} établies pour la catégorie des revêtementsétales d'un schémas donné, et dont une partie sont énoncés dans Exp V n\textsuperscript{\uline{o}}4 (conditions G1 à G3). De façon précise, on peut dire que dans la catégorie $\operatorname{Rev}^{\underline{R}}(S)$, les limites projectives finies et les limites inductives finies existent, et satisfont aux mêmes propriétés d'exactitude que dans la catégorie des ensembles (telle l'exactitude du foncteur changement de base dans cette catégorie); \deleted{de plus\unsure{;|,}} \inserted{en particulier} tout morphisme de la catégorie se factorise en le produit d'un épimorphisme et d'un mnomorphisme, \inserted{les épimorphismes et monomorphismes sont \deleted{\unsure{s}tri} effectifs; enfin} tout monomorphisme est un isomorphisme avec un sommande direct. On fera attention que le foncteur ``réalisation géométrique'' sur $\operatorname{Rev}^{\underline{R}}(S)$ commute bien aux limites inductives finies (et \deleted{par suite au\unsure{x}} en particulier aux sommes finies), mais non pas en général aux limites projectives finies, et en particulier, pas aux produits de deux objets.

Pour appliquer la théorie de Galois générale de V 4, il faut dis-

\scanpage{26}%
% archive sheet number: 25
poser de plus d'un \uline{foncteur fibre}
\[
F:\operatorname{Rev}^{\underline{R}}(S)\qquad(\mathrm{Ens\ finis}),
\]
i.e. d'un foncteur exact\inserted{(conditions G4,G5 de V 4)}. Dans le cas de l'exemple 1, tout point géométrique de $Z$ fournit un tel foncteur fibre, et ce foncteur est conservatif (condition G6) si on suppose que toute partie ouverte et fermée de $Z$ stable par $G$ est vide ou égale à $Z$. De même, dans le cas de l'exemple 2, on obtient un \deleted{\unsure{pai|par}} foncteur fibre en prenant un point géométrique de $Z=\lim Z_i$. Bour obtenir un foncteur fibre dans le cas plus général de l'exemple 3, on choisit un point géométrique $\qquad S$ de $S$ \deleted{tel que} qui n'appartienne pas à\deleted{\unsure{u|n}} ``l'ensemble de ramification''de $\underline{R}$, i.e. tel que pour tout $S'$ dans $\mathrm{Et}_S$ et tout objet $M$ de $\underline{C}_{S'}$, l\unsure{\deleted{'}e|\deleted{e}'} \deleted{objet} \inserted{revêtement à} \deleted{à} opérateurs $\underline{r}(M)$ \deleted{sur $S'$}\inserted{(groupe $G$)} de $S'$ soit principal en les points de $S'$au dessus du l'image \inserted{$s$} de $\quad$ dans $S$. Alors \deleted{pour tout revêtement étale} \deleted{$T$ à groupe $G$ \unsure{du|d} de \unsure{à'un|d'un|l'un} du schéma à opérateurs $G$} on constate que pour tout objet $T$ de $\operatorname{Rev}^{\underline{R}}(S)$, $\underline{r}(T)$ est étale enr $S$ en $s$, et le foncteur $T\quad\underline{r}(T)(\quad)$ (fibre géométrique en $\quad$) est bien exact. Pour qu'il soit conservatif, il faut et il suffit que le site $\underline{C}$ (pour la topologie image inverse de celle de $\underline{\mathrm{Et}}_S$) soit connexe; \deleted{i.e. que tout crible couvrant \unsure{du|de} ce \unsure{sjte|site}} si on suppose, avec les notations de l'exemple 3, que pour tout $M$, définissant $(\underline{r}(M),G)$, $\underline{r}(M)/G\qquad S'$ soit un isomorphisme (ce qui sera sans doute le cas dans tous les cas intéressants), alors il revient au même de dire que $S$ est connexe. Appliquant la théorie de V 4, on trouve donc un groupe pro-fini, noté ${}^{\underline{R}}_1(S,\quad)$, \deleted{qui classifi} et une équivalence de la catégorie des ensembles finis où ce groupe opère continûment, avec la catégorie $\operatorname{Rev}^{\underline{R}}(S)$.

3. \uline{Fonctorialités pour} $S$ \inserted{\uline{et} $\underline{R}$} \uline{variables}

\unsure{Crnsidérons|Considérons} un morphisme $f:T\qquad S$ de schémas, d'où un foncteur

\scanpage{29}
% archive sheet number: 28

est``l'image inverse'' de $\underline{R}_{(S)}$, toujours sous les conditions de l'exemple 3. Alors il est t\uline{rivial} que $\operatorname{Rev}^{\underline{R}_{(S)}}(S)\qquad\operatorname{Rev}^{\underline{R}_{(T)}}(T)$ est une équivalence, utilisant le résultat analogue pour les revêtements étales de schémas $X$ et $Xx_{S}T$ (où $X\ \deleted{\mathrm{est}}=\underline{s}(M)$ est un schéma \deleted{co\unsure{n|u}v\unsure{r}} sur $S$ ...). En particulier, quand on dispose d'un point géométrique non ramifié de $T$, on obtient un isomorphisme sur les groupes fondamentaux. Le même raisonnement, après localisation étale à peu près évidente, s'applique au\deleted{\unsure{g|q}} cas kummérien \deleted{a)}.

J'ai fait allusion p\unsure{\deleted{h}l|\deleted{l}h|\deleted{d}l|\deleted{l}d}usieurs fois à la donnée de définir une donnée de ramification image inverse d'une donnée de ramification sur $S$ par un morphisme $T\quad S$; la définition explicite serait nécessairement un peu lourde; je me borne à signaler que dans le cas où $\underline{R}$ est défini par $(Z,G)$ comme dans l'exemple 1, son image inverse est définie par $(Z_T=Zx_ST,G)$. L'ennui avec cette définition générale, c'est qu'elle ne s'applique pas directement au cas où $\underline{R}$ est une donnée kummérienne associée à un systême $(D_i)$, car la donnée image inverse au sens précis précédent n'\uline{est pas} la donnée correspondante au\deleted{$R$} systême $(D_{i\,T})$ même à équivalence de catégories fibrées près; elle lui est seulement ``essentiellement équivalente'' au sens du n$^{\underline{o}}$1. Seule une rédaction en forme permettra de se rendre compte s'il y a lieu d'introduire cette notion d'image inverse, ou s'il ne vaut pas mieux donner des conditions axiomatiques sur la donnée $\unsure{(*)|(\pi)|(\deleted{=}x)|(\deleted{x}=)}$, sous les conditions de l'exemple 3, qui assurent moralement que $\underline{R}_{(T)}$ est essentiellement équivalente à l'image inverse de $\underline{R}_{(S)}$. Bien sûr, ces conditions, pour \deleted{avoir} être utiles, doivent s'appliquer directement aux cas kummériens. Ce sera sans doute le parti le plus raisonnable. Il me semble également qu' serait raisonnable, assez rapidement après les définitions général-

\scanpage{31}
% archive sheet number: 30

passage au complété (par exemple si $\underline{O}_{S,s}$ à fibres formelles géom. normales, par exemple s'il est excellent). Les mêmes observations s\unsure{\deleted{d}'a|\deleted{'a}d}ppliqueront sans hypothèse de normalité ou de réduction quelconque; à condition de se borner aux revêtements \deleted{kummériens}\inserted{mod.ram.} \uline{galo\unsure{\deleted{s}i|\deleted{i}s|\deleted{n}i|\deleted{i}n}\unsure{\deleted{d}s|\deleted{s}d|\deleted{i}s|\deleted{s}i}iens} ; cf $n^{\underline{o}}$ suivant.

\unsure{\deleted{3}4|\deleted{4}3}. \uline{Cas des revêtements ga\unsure{\deleted{h}l|\deleted{l}h}oisiens}.

Un revêtement \uline{galoisien} de type $\underline{R}$ de $S$, à groupe de Galois $G$, est un objet $X$ de $\operatorname{Rev}^{\underline{R}}(S)$, sur lequel le groupe fini $G$ opère à \deleted{\unsure{r|n|x}}droite, de tel façon que dans la catégorie $\operatorname{Rev}^{\underline{R}}(S)$ on obtienne un\deleted{e} revêtement principal de l'objet final, i.e. $1^{\underline{o}}$) le morphisme $XxG\quad XxX$ est un isomorphisme et $2^{\underline{o}}$) le morphisme $X\quad\text{objet final}$ est couvrant. Comme $\operatorname{Rev}^{\underline{R}}(S)$ est en tous cas une catégorie multigaloisienne, la notion de revêtement galoisi\unsure{\deleted{i}e|\deleted{e}i|\deleted{t}e|\deleted{e}t}n de type $\underline{R}$ s'explicite également en termes des groupesfondamentaux décrivant cette catégorie. On peut enfin l'expliciter en revenant à la définition de $\operatorname{Rev}^{\underline{R}}(S)$, et en exigeant que les revêtements des $\underline{s}(M)$ $(M\quad\operatorname{Ob}\underline{C}_{S'}\ )$ correspondants à $X$ soient principaux de groupe $G$. Par abus de langage, l\unsure{\deleted{'}e|\deleted{e}'}\deleted{objet}\inserted{revêtement} à opérateurs $(\underline{r}(X),G)$ de $S$ s'appellera également revêtement galoisien de type $\underline{R}$ de $S$. Cet abus de langage ne présente pas d'inconvénient trop sérieux dans le cas où le foncteur de la catégorie des revêtements galoisiens de type $\underline{R}$ de $S$, dans la catégorie des revêtements à opérateurs de $S$, défini par le foncteur ``réalisation géométrique'', est pleinement fidèle,lorsqu'on se borne à prendre comme morphismes dans la première les seuls isomor\deleted{ph}phismes. (Alors, en effet, pour un revêtement à opérateurs $(T,G)$ donné de $S$, il existe, à isomorphisme canonique près, au plus un $(X,G)$ \deleted{dont} galoisien de type $\underline{R}$ dont il provient). Nous allons donner des conditi\deleted{o}ons pour qu'il en soit bien ainsi, en nous plaçant comme d'habitude dans le cas de l'exemple 3. Voici d'abord la condition néc. et suffisan

\scanpage{32}
% archive sheet number: 31

te dans le cas particulier de l'exemple 1: pour tout localise strict $S'$ d'un point de $S$, prenant une composante connexe $Z'_o$ de $Z'=Zx_{\unsure{\substack{S\\S}|\substack{\deleted{S}\\S}|\substack{S\\\deleted{S}}}}S'$ et son stabilisateur $H$ dans $G$, et pour deux sous-groupes \inserted{distingués} $H_1,H_2$ de $H$ et tout i\unsure{\deleted{s}d|\deleted{d}s}\unsure{\deleted{i}o|\deleted{o}i|\deleted{h}o|\deleted{o}h}morphisme $u:Z'_o/H_1^{\deleted{\unsure{1|I}}}\quad Z'_2/H_2$, \deleted{il \unsure{u|n}} compatible avec un isomorphisme $\handwritten{\phi}:\unsure{H|\deleted{G}H|\deleted{H}G|\deleted{\unsure{H}}H}/H_1\quad H/H_2$, il existe un élément $g$ de $\unsure{H|\deleted{G}H|\deleted{H}G|\deleted{\unsure{H}}H}$ tel que $\operatorname{int}(g)$ applique $H_1$ sur $H_2$ et induise $\handwritten{\phi}$ par passage au quotient, et tel que $u$ soit induit par l'endomorphisme $g_{Z'_o}$ de $Z'_o$ par passage aux quotients. Dans le cas général de l'exemple 3, il faut exiger ceci pour tous les couples $(\underline{s}(M),G)$ sur les divers $S'\quad\underline{\operatorname{Et}}_S$.

This condition is satisfied in the case of Kummer data, as one sees by combining 1.10. and 1.15. of Murre's notes. This phenomenon explains why I have been so reluctant (wrongly so) to define \deleted{Kummer}\inserted{tamely ramified} coverings except in the Galois case: in the non galois case, it will have to be not just a covering satisfying some extra \uline{condition\unsure{\deleted{\unsure{g|p}}s|\deleted{s}\unsure{g|p}|s|g|p}}, but a covering with some extra \uline{structure} ; whereas in the Galois case, the extra structure is uniquely determined, if it exists. Note however that \deleted{for morphis\unsure{m}} not all morphisms \inserted{(in the usual naive sense)} between (even Galois) tame coverings which are not isomorphisms should be considered, only those which come form morphisms in $\operatorname{Rev}^{\underline{R}}(S)$, - that one has to be careful is shown by the counterexample to prop.1.16. of Murre's notes. Also, the natural notion of products of tame \deleted{Galois} coverings, whether \unsure{\deleted{G}g|\deleted{g}G}alois or not, is not the naïve fibered product over $S$ (which by the way accounts for the complication of the proof of Murre's lemma 4.3.). This shows rather clearly that even in the Galois case, the consideration of the category $\operatorname{Rev}^{\underline{R}}(S)$ is undispensable for a good comprehension of what is really going on.

\scanpage{33}
% archive sheet number: 32

\uline{Rem\unsure{\deleted{r}o|\deleted{o}r}\unsure{\deleted{o}r|\deleted{r}o}ds} concernant l'exemple 3 du par.1. - La présentation de cet exemple n'est ni élégante ni naturelle, la donnée de $\underline{s}'$ et de \quad étant un élément de structure artificiel. Il vaut mieux présenter les choses ainsi, en termes de conditions sur $\underline{R}$, savoir les suivantes:

a) Pour tout objet $M$ d'un $\underline{C}_{S'}$, le faisceau \inserted{étale} associé au préfaisceau sur $\underline{\operatorname{Et}}_{S'}$ des automorphismes de $M$ est \deleted{d\unsure{u|e}} représentable par un schéma en groupes $\underline{G}(M)$ fini et étale sur $S$ (évidemment déterminé à isomorphisme unique près.

Il est évident que $G(M)$ opère sur le schéma relatif $\unsure{\underline{s}(M)=Z|\underline{s}(M)\equiv Z}$.

b) Pour tout objet $S'$ de $\underline{\operatorname{Et}}_S$, tout morphisme \inserted{\deleted{$M'\quad M$}} dans $\underline{C}_{S'}$ est un épimorphisme, et tout endomorphisme d'un obje\unsure{\deleted{y}t|\deleted{t}y} de $\underline{C}_{S'}$ est un automorphisme. Pour tout morphisme $\inserted{u:}M'\quad M$ dans $\underline{C}_{S'}$, et tout automorphisme $g'$ de $M'$, il existe un automorphisme $g$ de $M$ tel que $gu=ug'$.

Cet automorphisme est évidemment unique, de plus la condition précédente implique que à $\inserted{u:}M'\quad M$ est associé un unique homomorphisme d schémas en groupes $G(M')\quad G(M)$ tel que $u$ soit compatible avec les opér\unsure{a}tions de $G(M')$ et $G(M)$ sur $M'$ et $M$.

c) Avec les notations précédentes, \deleted{si} l'homomorphisme $G(M')\quad G(M)$ est surjectif, et si $H$ désigne son noyau, l\unsure{\deleted{'}e|\deleted{e}'}\deleted{homo}morphisme $\underline{s}(M')\quad\underline{s}(M)$ induit un \deleted{morphis\unsure{m}e} isomorphisme $\underline{s}(M')/H\quad\underline{s}(M)$, \deleted{i\unsure{.}e\unsure{.} fait de $\underline{s}(M')$ un \unsure{fibéé|fib\'{r}é} principal sur $\underline{s}(M)$, de groupe} (NB on suppose que les $\underline{s}(M)$ sont \deleted{\unsure{relativement} des \unsure{revêtem}} relativement affines, de sorte que le passage au quotient a un sens).

d) Deux homomorphismes de $M'$\deleted{d} dans $M$ sont,localement pour la top\unsure{\deleted{I}o|\deleted{t}o|\deleted{l}o|\deleted{o}I|\deleted{o}t|\deleted{o}l}logie étale, déduits l'un de l'autre par composition avec un automorphi\unsure{s} de $M$.

\scanpage{35}
% archive sheet number: 34

Dear Murre,\\

I am very sorry I did not succeed to convey the intuitive idea behin\inserted{d} the general nonsense of my notes. It seems to me that the basic example in order to understand the idea is example 1, where you can take $Z$ to be a standard Kummer covering for definiteness, $Z=Z_a^n$, and $S$ normal. Intuitively, when you look at coverings \inserted{(normal, say) $S'$} of $S$ whose ramification type is not worse than the one of $Z$ over $S$, you mean that the \inserted{normalize\unsure{\deleted{e}d|\deleted{d}e}} inverse image $Z'$ of $S'$ over $Z$ is étale. \deleted{B\unsure{u}t by transport de} From the birational point of view, assuming $S'$ connected and therefore corresponding to a field extension $K'$ of the field of functions $K$ of $S$, this means simply that $K'$ is isomorphic to a sub\inserted{-$K$-}extension of an extension \inserted{$L'$} of the function field $L$ of $Z$, unramified with respect to the model $Z$; when $S$ (and whence $Z$) is strictly local, \inserted{hence $L'=L$,} this means simply that $K'$ is isomor\unsure{\deleted{f}h|\deleted{h}f|\^{h}}ic to a subextension of $L$, a very reasonable way indeed of interpreting the intuitive fact that $S'/S$ has no worse ramification than $Z/S$. (NB We assume of course now,for simplicity, that $Z$ is connected, and mor\unsure{\.{e}|é|e\inserted{'}}over that $Z$ is the normalisation of $S$ in some finite Galois extension $L$ of $K$). Note that anyhow, we want that the notion in question should be of local nature for the étale topology, thus we can reduce always to the s rictly local case. Now by transport de structure the group $G$ operating on $Z$ must operate on $S'$\inserted{(in a way compatible with\deleted{h} operationson $Z$)}, and it is immediate that we recover $S'$ from the $G$-covering $Z'$ of $Z$ as being simply $S'=Z'/G$. Thus \deleted{\unsure{w|n}e re} the category of coverings of $S$ we are interested in is just equivalent to the catégory of \inserted{étale} coverings of $(Z,G)$. This latter catégory on the other hand is of a very standard type; the presence of $G$ does not prevent it to share all the properties of catégories of étale coverings

\scanpage{79}
Voorschoten, \handwritten{30-4-69.}\\
Palestrinalaan 11.\\

Dear Grothendieck,\\

Last fall I wrote you that I was planning to work out your notes on formal tamely ramified coverings together with a result on ``plumbing''.

In doing so it turned out to be necessary- at least \unsure{\deleted{i}b|\deleted{f}b|\deleted{r}b|\deleted{b}i|\deleted{b}f|\deleted{b}r}y the method \unsure{I|t} followed- to write again a lot of preliminaries on tamely ramified coverings, this time for formal schemes. \unsure{l|1|I} am still not ready with the above mentioned ``project''\handwritten{!} Since it took me already much time, I would appreciate it very much to have -before proceeding- your opinion about the \uline{preliminary version} which \unsure{I|1|l} have prepared so far. Of course it is not necessary to read the notes in detail. My question is: what do you think of this set-up and do you think it worthwhile to continue in this way? To me the entire thing seems somewhat indigestible.

I have adjoined some comments on the different chapters. The notes consists of two parts. Part one ( Chap. I-VII) are the preliminaries, part two ( CHap.VIII-X) are your notes and the plumbing. The chapter on the analogue of the theorem of Mumford has still to be written.

I hope you don't mind too much that I bother you again with these tamely ramified coverings\handwritten{!}

With best wishes,\\
Sincerely yours,\\
\handwritten{J.P. Murre}\\
% archive sheet number: 78

\scanpage{80}
\mspage{\handwritten{1\unsure{.}}}

\uline{Comments}.

\uline{Chap.I.} \'{E}tale morphisms\inserted{of formal schemes} are defined as inductive systems of étale morphisms of usual schemes (1.1.2 and EGA I 10.12.3).

Crucial lemmas: 1.1.6, 1.1.13 and 1.4.9.

As to 1.1.13: étale nbhs for formal schemes are locally completions of étale nbhs of usual schemes. This is used in 7.1.5, itself used in the reduction step 9.3.3.

\uline{Chap.II}. There are Kummer coverings (II-1), generalized Kummer coverings (II-2) and coverings of Kummer type (II-6)\handwritten{!}

No attempt has been made towards completeness; for instance there is no section on the automorphisms of (generalized) kummer coverings.

It would be useful to have a chara\inserted{\handwritten{c}}terization of the Algebra of a \uline{generalized} Kummer cov. (cf. with 2.1.8 for Kummer coverings).

\uline{Chap.III}. This is a nuisance\handwritten{!}

I have tried to write only down those things needed for Chap.VIII-X.

It was sufficient to follow the notes which \deleted{I} made for usual schemes and therefore I have not followed your new set-up. But one way or the other, one gets the feeling that the notion of tamely ramified covering, although natural, is very unmanageable if one has to use ``the functor'' instead of the geometric realization itself.

Lem\deleted{\unsure{o|c|e}}ma 3.2.11 is a \unsure{techni\deleted{c}val|techni\deleted{v}cal} lemma, used for instance in 3.5.8, 3.6.5 (and \deleted{many} at several places later on).

\uline{Very crucial are} \deleted{are} 3.4.1,3.4.2 ; also the notations introduced there are often used in the sequel.

For 3.5.5 I have given only an indication of the proof; in fact many proofs of III are standard but very unpleasant to spell out in \unsure{a\deleted{l}al|a\deleted{a}ll} details,

Part 3.6 (and especially 3.6.5) are important. Here the geometric realization determines the tamely ramified \unsure{\deleted{v}covering|\deleted{c}vovering}. Usually only normality, instead of regularity, seems to be needed. The difficulty is however that the notion of normality is (perhaps\unsure{'\_|)} not stable by étale localization (cf.1.7.3 and 1.7.4).

\uline{Chap.VI}. \deleted{\unsure{6|5|A}} 6.1.3 is used in8.1.\unsure{9}, see\inserted{\handwritten{page}} VIII-8 point c.

6.2.1 and 6.2.2 are technical results, used in the proof that 9.3.5 implies 9.3.2 (see \deleted{\unsure{see|IX|III|VII}} page IX-17).

% archive sheet number: 7\unsure{3|9}

\scanpage{81}
\mspage{(2).}

\uline{Comments}

\uline{Chap.VII}. This is a technical chapter. The results are delicate , but they are used in an essential way in 9.3 in the reduction step 9.3.3.

As pointed out before: 1.1.13 is used in 7.1.5. In 7.1.5 we have to work with Kummer coverings instead of generalized Kummer coverings; the reason (or at least one of the reasons) is that we don't have an analogous lemma for 2.1.8. As a consequence the assumption $D_i\handwritten{\cap}D_j=\handwritten{\varnothing}$ (see 9.0.1) has to be made.

Also 7.1.6 is a subtle lemma.

\uline{Chap.VIII}. The first page of your notes ``Revêtements formels étales... remark 8.1.6\unsure{.}\inserted{i.e. the isomorphism is \uline{canonical}} seems tome to be correct (\handwritten{!}) and essential for 8.1.9.iii.

\uline{Chap.IX}.

\unsure{Th\deleted{i}es|Th\deleted{e}is} is a working-out of yo\inserted{\handwritten{u}}r notes.

The assumption $D_{i,o}\handwritten{\cap}D_{j,o}=\handwritten{\varnothing}$ is-as p\unsure{\deleted{i}o|\deleted{o}i}inted out above- used in 9.3.3., in order to be able to apply 7.1.6 and 7.1.8 (page IX-13,line 4 top).

\uline{In IX-1} there are several points about which I am not certain (for instance page IX-4 below). Nevertheless I have used them (\handwritten{!}) because it seems to me that you used them implicitely and without mentioning in your notes and probably these things are more or less standard results in étale cohomology.

\uline{In IX-2} I have not included (yet\unsure{)|'} the proof of the key-theorem 9.2.3.

(this proof is explicit\unsure{\deleted{e}\handwritten{ly}|\handwritten{ly}} given in the letter of Giraud).

\uline{As to IX-3} : the reduction step 9.3.3. seems to me to be very delicate and depends upon the results in chap.VII. (In your notes you give the motivation in the margin \handwritten{!}). It is -in particular- at this point that one has the feeling that the notion of tamely ramified covering is cumbersome.

I hope that part IX-3 is correctly written\inserted{\handwritten{and that I don't have overlooked difficulties}}; I had a lot of trouble with it.

One final remark: I needed here the results on tamely ramified coverings for \uline{usual schemes}; are th\unsure{\deleted{o}e|\deleted{e}o}se included in your planned final version of SGA 1 Exp.XII (see 9.0.4,9.1,4, page IX-6 below)\handwritten{?}

\uline{Chap.X} The plumbing. X-3 is again very cumbersome. It \unsure{\deleted{o}w|\deleted{s}w|\deleted{w}o|\deleted{w}s}ould be nicer to work directly with the geometric realization, but there I encountered difficulties, therefore I followed an easy but long way.

% archive sheet number: 80
